\section{Introduction}

Data privacy is a critical concern in distributed applications
including privacy-preserving machine learning
\cite{li2021privacy,knott2021crypten,koch2020privacy,liu2020privacy}
and blockchains
\cite{ishai2009zero,lu2019honeybadgermpc,gao2022symmeproof,tomaz2020preserving}.
Secure multiparty computation (MPC) is a key software-based enabling
technology for these applications. MPC protocols provide both confidentiality
and integrity properties, formulated as real/ideal (aka simulator)
security and universal composability (UC) \cite{evans2018pragmatic},
with well-studied manual proof methods \cite{Lindell2017}.  MPC
security semantics has also been reformulated as
\emph{hyperproperties}
\cite{8429300,10.1145/3453483.3454074,skalka-near-ppdp24} to support
automated language-based proof methods. Recent research in the SMT
community has developed theories of finite fields \cite{SMFF} with
clear relevance to verification of cryptographic schemes that rely on
field arithmetic. The goal of this paper is to combine hyperproperty
formulations of security with SMT verification as a foundation for
enforcing correctness and security properties of MPC protocols.

Our focus is on low-level MPC protocols. Previous high-level
MPC-enabled languages such as Wysteria \cite{rastogi2014wysteria} and
Viaduct \cite{10.1145/3453483.3454074} are designed to provide
effective programming of full applications, and incorporate
sophisticated compilation techniques such as orchestration
\cite{viaduct-UC} to guarantee high-level security properties. But
these approaches rely on \emph{libraries} of low-level MPC protocols,
such as binary and arithmetic circuits. These low-level protocols are
probabilistic, encapsulate abstractions such as secret sharing and
semi-homomorphic encryption, and are complementary to high-level
language design.

\subsection{The $\metaprot/\minicat$ Language Framework}

Our language framework incorporates two languages. The $\minicat$
language supports direct definition of message passing protocols.  In
this work we develop a \emph{share type} system enforcing the
confidentiality property of MPC which is called passive simulator
security. We also develop and integrity type systems for $\minicat$
that enforces malicious MPC security. Both share types and integrity
types are supported with the use of an SMT-style constraint system.
In addition to security properties, the constraint system also
supports protocol correctness properties.

While $\minicat$ provides a foundation for low-level MPC protocol
specification, it does not provide higher-level features such as
function definitions, and the share and integrity type systems are not
automated. So we define the $\metaprot$ language as a higher-level
metaprogramming language that evaluates to $\minicat$ protocols.
$\metaprot$ supports compositional definition of protocol components
and also supports automation of constraint verification and type
checking especially by leveraging SMT automation.  Thus $\metaprot$
program analysis automates security and correctness properties of the
low-level $\minicat$ protocols that $\metaprot$ programs evaluate to.

The specification of an MPC protocol is called its \emph{ideal
functionality} and a standard approach \cite{269581} is to implement
ideal functionalities as either arithmetic or boolean circuits. An
advantage of this is that subprotocols for primitive operations, such
as addition and multiplication, or logic gates, can be implemented and
optimized as subprotocols. Consider the example of a two-bit adder
binary circuit in Figure \ref{fig-adder}. The ideal functionality of
the circuit is two-bit vector addition (ignoring carry). On the left
is a $\metaprot$ program that implements the protocol by ``wiring''
secure logic gates implemented using the Goldreich-Micali-Wigderson
(GMW) protocol~\cite{goldreich1987play}. These are encapsulated in the $\ttt{xorgmw}$ and
$\ttt{andgmw}$ definitions. In the first lines of the program, clients
1 and 2 each generate secret shares of their secret bits-- $x_1$ and
$x_2$ for client 1, and $y_1$ and $y_2$ for client 2. Shares are
identified in the program by strings $\ttt{"x1"}$, $\ttt{"y1"}$, etc.,
as are the gates in the protocol. So the circuit is wired in the
same way as the logic circuit diagram on the right. Each gate
specifies a subprotocol where the operations are performed on
encrypted wire data-- in the case of GMW, these are secret shares. 


\begin{fpfig}[t]{A two-bit adder circuit (without carry) in $\metaprot$ with
    verification postcondition, using the GMW protocol.}{fig-adder}
\begin{center}
\begin{tabular}{cc}
\begin{tabular}[b]{l}
\verb+main() {+\\
\verb+   genshares("x1", 1, 2); genshares("x2", 1, 2);+\\ 
\verb+   genshares("y1", 2, 1); genshares("y2", 2, 1);+\\ 
\verb+   xorgmw("z1", "x1", "y1");+ \\
\verb+   andgmw("g1", "x1", "y1");+\\
\verb+   xorgmw("g2", "x2", "y2");+\\
\verb+   xorgmw("z2", "g2", "g1")+\\
\verb+}+ \\
\verb+postcondition:+\\
\verb+  (|RECON("z2"), RECON("z1")| ==+\\
\verb+      BVAdd(|RECON("x2"), RECON("x1")|,+\\
\verb+            |RECON("y2"), RECON("y1")|))+
\end{tabular}
& 
{
\scriptsize
\begin{circuitikz}
  \draw
  (0,0)
  node[label=left:$x_1$] {}
  to[short, o-]
  (1,0)
  node[xor port, anchor=in 1] (z1) {}
  (0,52 |- z1.in 2)
  node[label=left:$y_1$] {}
  to[short, o-]
  (z1.in 2)
  (z1.out)
  to[short, -o]
  ++(3,0)
  node[label=right:$z_1$] {}
  (2,-1.5)
  node[and port, anchor=in 1] (g1) {}
  (g1.in 1)
  -- ++(-1,0)
  to[short,-*] (1,0)
  (g1.in 2)
  -- ++(-1.5,0)
  to[short, -*] (.5, 52 |- z1.in 2)
  (0,-3)   
  node[label=left:$x_2$] {}
  to[short, o-]
  (1,-3)
  node[xor port, anchor=in 1] (g2) {}
  (0,52 |- g2.in 2)
  node[label=left:$y_2$] {}
  to[short, o-]
  (g2.in 2)
  (g1.out)
  to[short,-]
  ++(.2,0)
  node[xor port, anchor=in 1] (z2) {}
  (g2.out)
  -- ++(1.2,0)
  to[short, -] (3.75, 52 |- z2.in 2)
  (z2.out)
  to[short, -o]
  ++(.25,0)
  node[label=right:$z_2$] {}  
  ;
\end{circuitikz}
}
\end{tabular}
\end{center}
\end{fpfig}


An important feature of the example is the $\ttt{postcondition}$ that
specifies the correctness of the circuit. Here the $\ttt{|...|}$
syntax denotes a concatenation of bits into a bit vector, and
$\ttt{BVAdd}$ is bit vector addition. These operations are components
of the SMT cvc5 Bit-Vector theory and the postcondition is
automatically verified using it. $\ttt{RECON}$ is a constraint-level
function indicating the reconstructed value represented by secret
shares. So the postcondition says that the reconstructed shares of
the outputs are equal to the sum of the reconstructed secret shares
of the inputs.

In the remainder of the paper we will go into more complete detail
about the underyling definitions of $\ttt{genshares}$, $\ttt{xorgmw}$,
and $\ttt{andgmw}$ (see in particular Figure \ref{fig-gmwprepost} in
Section \ref{section-metalang} and Figure \ref{fig-gmwsharetypes} in
Section \ref{section-metaty}). This includes their share type checking that
guarantees passive simulator security in this and other examples.
Another important point is that other protocols can be used to define
gates, for example malicious secure protocols, since $\metaprot$
and $\minicat$ support low-level protocol specification.

\subsection{A Brief Introduction to MPC Protocols}

An MPC protocol allows a collection of parties to compute a function
of their private inputs without revealing those inputs to each other.
The function to be computed is typically called the \emph{ideal
  functionality}, and a correct protocol must produce the same output
as the ideal functionality. An MPC protocol is secure if it reveals no
additional information to the parties beyond what is intentionally
revealed by the ideal functionality. The information available to a
party (its input, randomness, received messages, and outputs) is
called its \emph{view}; an adversary can observe the views of all
corrupted parties~\cite{evans2018pragmatic,Lindell2017}.

\emph{Secret sharing} splits a secret value into \emph{shares} that
allow limited computations but prevent individual parties from
learning the secret. \emph{Additive secret sharing} is one example: to
additively share $x \in \mathbb{F}_p$ between two parties, we sample a
uniform $r \in \mathbb{F}_p$ and gives the shares $x_1=r$ and
$x_2=x-r$ to the respective parties. To reconstruct the secret, we
compute $x_1+x_2=x$. Each share by itself is uniform and therefore
does not reveal $x$, but two shares together uniquely determine $x$.
This construction is a finite-field \emph{one-time pad} (OTP): for
every possible $x$, the observed value $x-r$ has the same uniform
distribution, so even an adversary with unlimited computational power
learns nothing about $x$. The mask must be fresh for every secret. For
example, reusing $r$ to mask $x$ and $y$ reveals their difference,
since $(x-r)-(y-r)=x-y$. In this paper $\ttt{RECON}$ denotes
reconstruction, and over $\mathbb{F}_2$ additive reconstruction is
XOR, as used for GMW wire values \cite{goldreich1987play, evans2018pragmatic}.

Additive secret sharing is a form of \emph{semi-homomorphic}
encryption: it allows limited forms of computation without decryption.
Adding two additive sharings componentwise gives a sharing of the sum.
However, additive secret shares cannot be multiplied to obtain shares
of the product; to support multiplication, MPC protocols use
additional machinery (e.g. oblivious transfer in GMW, or preprocessed
Beaver triples in arithmetic protocols~\cite{goldreich1987play,
  evans2018pragmatic, BDOZ,10.1007/978-3-030-68869-1_3}).

\emph{Oblivious transfer} (OT) is a two-party protocol that allows one
party to select a value from the other without revealing which one was
selected. The \emph{sender} supplies two values $m_0,m_1$, and the
\emph{receiver} supplies a choice bit $b$. The receiver obtains $m_b$
without learning $m_{1-b}$, while the sender learns nothing about $b$.
For example, GMW uses OT to select an entry from a table that
implements a Boolean gate without revealing either party's
secret-shared input bits. Different OT implementations rely on
different assumptions, so our language treats OT abstractly and
requires only these functional and privacy guarantees
\cite{goldreich1987play, evans2018pragmatic}.

Another approach moves expensive work into an input-independent
\emph{preprocessing} phase that distributes correlated randomness
before the parties know their private inputs. A Beaver
triple~\cite{beaver1991efficient} consists of secret shares of random
values $a,b,c$ satisfying $c=ab$. To multiply shared $x$ and $y$, the
parties reconstruct the masked differences $d=x-a$ and $e=y-b$, then
locally obtain shares of $xy=c+db+ea+de$. Reusing a Beaver triple
leads to correlations that break security, so triples and other
preprocessing resources must be fresh where required and securely
generated; this paper specifies their needed properties without
implementing that generation protocol \cite{evans2018pragmatic,
  BDOZ,10.1007/978-3-030-68869-1_3}.

MPC confidentiality is conventionally defined by comparing a
\emph{real} protocol execution with an \emph{ideal} execution in which
a trusted party receives the inputs, evaluates the ideal
functionality, and returns only its specified outputs. A simulator for
a corrupt coalition receives only the information available in the
ideal execution---in particular, the coalition's inputs and prescribed
outputs---and must generate a simulated view indistinguishable from
its real protocol view. The existence of such a simulator means that
\emph{anything learned from the real messages could also have been
  learned from the allowed ideal information}---in other words, the
real protocol reveals no more than the ideal
functionality~\cite{Lindell2017}.

Depending on the setting, indistinguishability may mean equality of
distributions, statistical closeness, or computational
indistinguishability. Equality gives \emph{perfect}, or
\emph{information-theoretic}, security even against unlimited
computation. Statistical security permits a small difference between
the distributions, bounding the distinguishing advantage even for
unlimited computation. Computational security instead considers
efficient adversaries and rests on a hardness assumption. Its
\emph{security parameter}, such as a key length (or a field size for
field-based authentication), controls both cost and protection; a
distinguishing advantage is \emph{negligible} when it decreases faster
than any inverse polynomial as that parameter grows. This paper uses
equality of finite distributions in its core model.

The adversary model determines how the corrupt parties behave. A
\emph{passive} (or honest-but-curious) adversary records and analyzes
the combined view of corrupt parties, but those parties follow the
protocol. A \emph{malicious} (or active) adversary may instead choose
arbitrary inputs, messages, and local computations. Malicious security
therefore combines confidentiality and
integrity~\cite{evans2018pragmatic}.

BDOZ and SPDZ protect integrity using \emph{authenticated
  shares}~\cite{SPDZ1, SPDZ2, BDOZ, 10.1007/978-3-030-68869-1_3}. A
\emph{message authentication code} (MAC) is a tag that binds a value
to secret authentication information; in these protocols, another
party holds the information needed to check a share's tag. A corrupt
party cannot tamper with a share, since it is impossible to produce a
corresponding valid tag. The MACs are linearly homomorphic, so tags
can be updated consistently when shares are added. Before revealing an
authenticated value, the parties check the resulting tag and abort if
it is invalid.

Ideal outputs may \emph{intentionally} reveal information, so MPC
security differs from ordinary noninterference. MPC security requires
only that the protocol's execution reveal nothing \emph{additional}
compared to the ideal functionality. Section \ref{section-model}
formalizes this idea as \emph{noninterference modulo output} and uses
it as a hyperproperty sound for passive simulator
security.

\subsection{Summary of the Theory and Implementation}

Here we summarize the main theoretical developments and results in
the paper and their connections.  In Section \ref{section-lang}
we define the $\minicat$ language model. In Section \ref{section-model}
we articulate our security model as hyperproperties of
noninterference modulo output (Definition \ref{definition-NIMO})
and integrity (Definition \ref{def-integrity}). As has been
shown in previous work \cite{skalka-near-ppdp24, skalka-near-esop25},
noninterference modulo output is sound for passive simulator
security, while the addition of integrity obtains malicious security.
Figure \ref{fig-framework-workflow} summarizes how the languages,
analyses, and guarantees fit together.

\begin{figure*}[t]
\centering
\resizebox{\textwidth}{!}{%
\begin{tikzpicture}[
  font=\sffamily\footnotesize,
  language/.style={
    draw=blue!60!black, fill=blue!8, rounded corners=2pt,
    align=center, text width=27mm, minimum height=9mm, inner sep=3pt
  },
  analysis/.style={
    draw=orange!70!black, fill=orange!10, rounded corners=2pt,
    align=center, text width=28mm, minimum height=9mm, inner sep=3pt
  },
  guarantee/.style={
    draw=green!50!black, fill=green!9, rounded corners=2pt,
    align=center, text width=28mm, minimum height=9mm, inner sep=3pt
  },
  implementation/.style={
    draw=black!55, fill=black!4, rounded corners=2pt,
    align=center, text width=27mm, minimum height=9mm, inner sep=3pt
  },
  group/.style={
    draw=black!45, dashed, rounded corners=3pt, inner sep=6pt
  },
  flow/.style={
    -{Latex[length=2mm]}, thick, draw=black!70
  },
  link/.style={
    thick, draw=black!70
  },
  edge label/.style={
    font=\scriptsize, inner sep=0pt, yshift=4pt, text=black!80
  }
]

\node[language] (prelude) at (0,4.1)
  {\textbf{Prelude}\\language and annotations\\Sec.~\ref{section-metalang}};
\node[analysis] (prelude-hoare) at (0,2.4)
  {Automated abstract\\Hoare logic\\Sec.~\ref{section-metahtrip}};
\node[analysis] (prelude-types) at (0,0.7)
  {Automated share\\type checking\\Sec.~\ref{sec:metapr-type-check}};

\node[language] (overture) at (4.2,4.1)
  {\textbf{Overture}\\protocol semantics\\Sec.~\ref{section-lang}};
\node[analysis] (overture-hoare) at (4.2,2.4)
  {SMT interpretation\\and Hoare logic\\Sec.~\ref{section-smt}};
\node[analysis] (share-types) at (4.2,0.7)
  {Share types\\Sec.~\ref{section-sty}};
\node[analysis] (integrity-types) at (4.2,-1.0)
  {Integrity types\\Sec.~\ref{section-ipj}};

\node[guarantee] (correctness) at (8.4,2.4)
  {Functional correctness\\of pre/postconditions};
\node[guarantee] (passive) at (8.4,0.7)
  {Noninterference modulo output\\and passive security\\Sec.~\ref{section-model}};
\node[guarantee] (integrity) at (8.4,-1.0)
  {Integrity\\Sec.~\ref{section-model}};
\node[guarantee] (malicious) at (12.3,-0.15)
  {Malicious security\\Sec.~\ref{section-model}};

\node[group, fit=(prelude)(prelude-hoare)(prelude-types),
  label={[font=\bfseries\small,text=black!70]above:Prelude language}] {};
\node[group, fit=(overture)(overture-hoare)(share-types)(integrity-types),
  label={[font=\bfseries\small,text=black!70]above:Overture language}] {};
\node[group, fit=(correctness)(passive)(integrity)(malicious),
  label={[font=\bfseries\small,text=black!70]above:Guarantees}] {};

\draw[flow] (prelude) -- node[edge label] {evaluates to} (overture);
\draw[flow] (prelude-hoare) -- node[edge label] {instantiates} (overture-hoare);
\draw[flow] (prelude-types) -- node[edge label] {preserves} (share-types);
\draw[flow] (overture-hoare) -- node[edge label] {proves} (correctness);
\draw[flow] (share-types) -- node[edge label] {enforces} (passive);
\draw[flow] (integrity-types) -- node[edge label] {enforces} (integrity);

\coordinate (security-join) at (10.35,-0.15);
\draw[link] (passive.east) -| (security-join);
\draw[link] (integrity.east) -| (security-join);
\draw[flow] (security-join) -- node[edge label] {} (malicious.west);

% \node[implementation, text width=23mm] (source) at (-0.5,-3.4)
%   {Prelude\\source};
% \node[implementation, text width=21mm] (parser) at (2.25,-3.4)
%   {ANTLR4\\parser};
% \node[implementation, text width=34mm] (backend) at (5.55,-3.4)
%   {Scala AST and backend\\compositional Hoare analysis};
% \node[implementation, text width=32mm] (solver) at (9.45,-3.4)
%   {cvc5 Java bindings\\Bit-Vector or finite-field solver};
% \node[implementation, text width=29mm] (result) at (13.0,-3.4)
%   {verified pre/postconditions\\and experiments\\Sec.~\ref{section-examples}};

% \node[group, fit=(source)(parser)(backend)(solver)(result),
%   label={[font=\bfseries\small,text=black!70]above:
%     implemented JAR workflow (Sec.~\ref{section-examples})}] {};

% \draw[flow] (source) -- (parser);
% \draw[flow] (parser) -- (backend);
% \draw[flow] (backend) -- (solver);
% \draw[flow] (solver) -- (result);

\end{tikzpicture}
}
\caption{Relationships among the paper's formal components and security
guarantees. Blue boxes denote languages and their semantics, orange
boxes denote verification analyses, and green boxes denote the
correctness and security guarantees they establish. Dashed outlines
group components by abstraction level.}
\label{fig-framework-workflow}
\end{figure*}

In Section \ref{section-smt} we develop an SMT-based Hoare logic
endowed with notions of validity and derivability, where the latter is
shown to imply the former (Lemma \ref{lemma-hoare-sound}). This is an
important foundation both for our verification method and for notions
of constraint correctness in type systems.  In Section
\ref{section-sty} we develop a novel share type system, enabling
security verification for protocols that use secret sharing as a
building block; our approach works for any secret sharing scheme and
enforces noninterference modulo output and therefore passive security
(Theorem \ref{theorem-sty}). The share type system also supports
verification of protocols with respect to ideal functionality.  In
Section \ref{section-ipj}, we develop a different type system that
enforces integrity (Theorem \ref{theorem-ipj}). Thus both share and
integrity typability of a $\minicat$ protocol imply that it is
malicious secure.

In Section \ref{section-metalang} we define the $\metaprot$ syntax and
semantics, including pre- and postcondition annotations on functions
(as in the example in Figure \ref{fig-adder}). Intuitively, any
$\metaprot$ program evaluates to a $\minicat$ protocol, so we refer to
it as a metaprogramming language. We define an automated Hoare logic
for $\minicat$ programs, to verify pre- and postcondition
annotations. This automated Hoare logic is shown to be sound in the
sense that pre- and postconditions will be valid for any function
applications (Lemma \ref{lemma-metahtrip}). An implication of this
(Corollary \ref{cor-metahtrip}) means that function pre- and
postconditions are compositional and need only be verified once,
achieving greater performance efficiency. In Section
\ref{sec:metapr-type-check}, we develop an automated type system for
$\metaprot$ that guarantees that any well-typed $\metaprot$ program
evaluates to a share-typable $\minicat$ protocol (Theorem
\ref{theorem-metaty}) and hence its passive security.

In Section \ref{section-examples} we develop and discuss extended
example applications of our type analyses to real protocols including
the Goldreich-Micali-Wigderson (GMW), Ben-Or-Goldwasser-Widgerson
(BGW), and Bendlin-Damgard-Orland-Zakarias (BDOZ). We also summarize
our implementation of the $\metaprot$ Hoare logic and evaluate
verification of both binary and arithmetic gate and circuit
properties. Our implementation integrates with cvc5 SMT to support
efficient verification, including Satisfiability Modulo Finite Fields
(SMFF) \cite{SMFF} for binary and arithmetic circuits, and the
Bit-Vector theory for binary circuits. The $\metaprot$ Hoare logic
analysis has been fully implemented in the current work in an
Antlr/Scala/cvc5-based workflow.

The implementation that produced our experimental results is available in our
companion
repository\footnote{\url{https://github.com/uvm-plaid/prelude-public}}. It uses
ANRLR4 to parse Prelude the source code, Scala to convert the Prelude abstract
syntax tree into cvc5 constraints, and the cvc5 solver to prove that for all
functions the precondition and function body meet the postcondition. When
built, the tool runs as a single Java archive (JAR). The repository includes
installation instructions, source code, further documentation, and a variety of
Prelude examples.


\subsection{Related Work and Our Contributions}
\label{section-related-work}

\compwrapfig

The goal of the $\metaprot/\minicat$ framework is to automate
correctness and security verification of low-level protocols.  Other
prior work has considered automated verification of high-level
protocols and manual verification of low-level protocols.  We
summarize this comparison in Figure \ref{fig-comp-wrap}.  In this
Figure we consider related systems in several dimensions.  This
inclues whether the language system provides \emph{probabilistic
features} (e.g., sampling) that are fundamental to low-level
protocols. It also includes whether existing verification methods
include \emph{probabilistic conditioning} that is essential to
condition analysis of information leakage inherent in ideal
functionalities. We consider whether the system is designed for
\emph{low-level} protocols, as opposed to being a framework for
integrating them in applications (e.g., during compilation). We
consider whether the system has integrated methods and/or mechanisms
for proving \emph{passive} and \emph{malicious security}, and whether
sound \emph{hyperproperties} have been formulated to support
enforcement of security via program analysis such as type systems.
Finally, we consider whether proof/enforcement mechanisms are
\emph{automated} in the system, vs., e.g., realized in proof
assistants.

The most closely related works are previous versions of the
$\metaprot/\minicat$ framework
\cite{skalka-near-ppdp24,skalka-near-esop25}. In the earliest work
\cite{skalka-near-ppdp24} foundations for $\metaprot/\minicat$ are
introduced including a verification algorithm. However, this algorithm
only works for $\mathbb{F}_2$-- e.g., binary circuits-- and the
datalog-based automation is a brute-force whole program analysis.
Subsequent work \cite{skalka-near-esop25} extends the language model
to arbitrary prime fields and introduces the use of SMT to optimize
verification in prime fields, though the automation is still a
whole-program analyis. The type systems presented in these works
support a gradual release property, weaker than noninterference
modulo output, and integrity properties. The current work provides
several advancements of $\metaprot/\minicat$:
\begin{itemize}
\item We develop a compositional Hoare logic that provides a foundation for
  compositional verification and type safety results.
\item We define a share type system that enforces noninterference modulo
  output as a result of type safety. This implies that share
  types enforce passive security and, together with integrity
  typability, imply malicious security.
\item Share types also support verification of ideal functionality
  specification of protocol outputs.
\item We extend the share type system to $\metaprot$, along with
  a compositional language of pre- and postconditions for functions.
  Type safety in $\metaprot$ ensures that share typing in programs
  is preserved by evaluation to $\minicat$ protocols.
\item The extensions in this work additionally enable the type-based
enforcement of noninterference modulo output for more complex secret
sharing protocols, e.g., that leverage Beaver triples.
\item We implement and evaluate the compositional verification of
  pre- and postconditions by considering examples of large circuits,
  specifically Bristol circuits. This evaluation demonstrates
  empirical advantes of compositional vs. whole-program verification.
\end{itemize}

The SAW/Cryptol framework \cite{10.1007/978-3-319-48869-1_5} supports
definitions of cryptographic protocols and verification using SMT.
The Swanky extension \cite{swanky} is especially related to our work,
since it is intended to apply to MPC protocols in binary and prime
fields. For example, Swanky has been used in the implementation of ZK
systems \cite{cryptoeprint:2020/1410}. However, while the goals of our
system are overlapping with regard to verification of protocol
correctness, we are also focused on type systems for enforcement of
security properties, incorporating reasoning about probabilistic
behavior of protocols. Our language model also more readily
accommodates primitives for message passing between multiple parties.
But as we discuss in Section \ref{section-conclusion}, leveraging
Swanky to accelerate verification is an interesting direction for
future work.

Other low-level languages with probabilistic features for privacy
preserving protocols have been proposed. The
$\lambda_{\mathrm{obliv}}$ language \cite{darais2019language} uses a
type system to automatically enforce so-called probabilistic trace
obliviousness.  But similar to previous work on oblivious data
structures \cite{10.1145/3498713}, obliviousness is related to pure
noninterference, not the conditional form related to passive MPC
security (Definition \ref{definition-NIMO}). A Haskell-based security
type system has been defined as an embedded DSL \cite{6266151} that
enforces a version of noninterference that is sound for passive
security, but does not consider malicious security. Properties of
real/ideal passive and malicious security for a probabilistic language
have been (manually) formulated in EasyCrypt \cite{8429300}.

Several high-level languages have been developed for writing MPC
applications. Previous work on analysis for the
SecreC language \cite{almeida2018enforcing,10.1145/2637113.2637119} is
concerned with properties of complex MPC circuits, in particular a
user-friendly specification and automated enforcement of
declassification bounds in programs that use MPC in subprograms. The
Wys$^\star$ language \cite{wysstar}, based on Wysteria
\cite{rastogi2014wysteria}, has similar goals and includes a
trace-based semantics for reasoning about the interactions of MPC
protocols. Their compiler also guarantees that underlying
multi-threaded protocols enforce the single-threaded source language
semantics. These two lines of work were focused on passive
security. The Viaduct language \cite{10.1145/3453483.3454074} has a
well-developed information flow type system that automatically
enforces both confidentiality and integrity through hyperproperties
such as robust declassification, in addition to rigorous compilation
guarantees through orchestration \cite{viaduct-UC}. However, these
high level languages lack probabilistic features and other
abstractions of low-level protocols, the implementation and security
of which are typically assumed as a selection of library components.

Program logics for probabilistic languages and specifically reasoning
about properties such as joint probabilistic independence is also
important related work. Probabilistic Separation Logic (PSL)
\cite{barthe2019probabilistic} develops a logical framework for
reasoning about probabilistic independence (aka separation) in
programs, and they consider several (hyper)properties, such as perfect
secrecy of one-time-pads and indistinguishability in secret sharing,
that are critical to MPC. Lilac \cite{li2023lilac} extends this line
of work with formalisms for conditional independence which is 
also particularly important for MPC \cite{skalka-near-ppdp24}.
